\doublespacing % Do not change - required

\chapter{Background}
\label{ch2}

%%%%%%%%%%%%%%%%%%%%%%%%%%%%%%%%%%%%%%%
% IMPORTANT
\begin{spacing}{1}
\minitoc
\end{spacing}
\thesisspacing
%%%%%%%%%%%%%%%%%%%%%%%%%%%%%%%%%%%%%%%

\begin{bodytwocol}

\section{Distributed and Federated Learning}
Distributed learning studies collaborative optimisation across multiple workers coordinated by a server: the server broadcasts a global model, workers compute local messages (gradients, momenta, or model deltas) on private data, and an aggregation rule forms the next global iterate. Federated learning (FL) is a widely adopted instance of this paradigm, motivated by communication efficiency and privacy considerations, where only model updates are exchanged rather than raw data (FedAvg and model averaging in \cite{DBLP:journals/corr/McMahanMRA16,pmlr-v54-mcmahan17a}). A recurring message in FL surveys is that real deployments are shaped by both system constraints (communication limits, stragglers, partial participation) and statistical heterogeneity (non-IID and unbalanced client data) \cite{KairouzMcMahanAvent:21}, which makes vanilla distributed optimisation assumptions fragile. Accordingly, a broad line of work refines federated optimisation under heterogeneity, including drift-control via control variates (SCAFFOLD) \cite{DBLP:journals/corr/abs-1910-06378} and adaptive server-side optimisation (FedOpt/FedAdam-style) \cite{reddi2021adaptive}. These methods primarily improve benign training stability and speed, but they do not directly address adversarial manipulation of worker messages, motivating a robustness literature whose conclusions must be interpreted jointly with heterogeneity and deployability constraints \cite{KairouzMcMahanAvent:21}.

\section{Byzantine Threats}
Robustness in distributed learning is commonly formalised using the Byzantine threat model originating from the Byzantine Generals problem \cite{10.1145/357172.357176}, where an unknown subset of workers may behave arbitrarily and possibly coordinate, sending corrupted or strategically crafted messages to the server. In FL, Byzantine behaviour can correspond to compromised devices, malicious participants, or severe faults; since the server typically only observes communicated messages, the aggregation rule becomes the critical security boundary. A complementary and often more protocol-realistic view categorises adversarial behaviour as poisoning: \emph{model poisoning} crafts malicious updates directly, while \emph{data poisoning} corrupts local data so that protocol-compliant training yields harmful but plausible messages \cite{KairouzMcMahanAvent:21}. This distinction matters because worst-case Byzantine models allow unbounded disturbances, whereas structured poisoning can induce bounded deviations that interact strongly with benign heterogeneity and stochastic noise.

\section{Robust Aggregation}
Under Byzantine/model-poisoning threats, most defenses design robust aggregators intended to limit the influence of poisoned messages so the aggregated direction remains close to the direction induced by regular workers. Representative mechanisms include selection-based rules such as Krum \cite{NIPS2017_f4b9ec30}; coordinate-wise robust statistics such as coordinate-wise median and trimmed mean \cite{pmlr-v80-yin18a}; and high-dimensional robust estimation perspectives for FL, including geometric-median-type estimators \cite{Pillutla_2022}, as well as broader Byzantine gradient-descent formulations \cite{10.1145/3154503}. However, robust aggregation is not solved by a single rule: El Mhamdi \emph{et al.} \cite{pmlr-v80-mhamdi18a} show that selection-based Byzantine-resilient learning can exhibit hidden vulnerabilities under coordinated adversaries, motivating threat-aware evaluation. Moreover, when benign update variance is high, attackers can craft perturbations that remain statistically plausible yet still cause large damage; Baruch \emph{et al.} \cite{BarBarGol:19} demonstrate that this variance shelter can enable stealthy manipulation that evades several defenses.

\section{Poisoning Attacks}
Poisoning attacks in FL span targeted vs untargeted objectives and can be instantiated through either model manipulation or data manipulation \cite{KairouzMcMahanAvent:21}. Backdoor attacks illustrate targeted poisoning: attackers preserve benign accuracy while inducing trigger-activated misbehaviour, often by amplifying their contribution via model replacement or scaling so that simple averaging is dominated by the malicious update \cite{DBLP:journals/corr/abs-1807-00459}. More generally, local model poisoning can be designed to explicitly degrade learning while evading Byzantine-robust aggregators \cite{10.5555/3489212.3489304}. Beyond magnitude-based scaling, Model Shuffle Attack (MSA) \cite{YANG2023158} shuffles parameter locations and applies scaling to degrade training in differential-privacy-based FL pipelines while remaining difficult to detect by naive anomaly checks.

\section{Label Poisoning}
While classical Byzantine attacks allow arbitrary messages \cite{10.1145/357172.357176}, many realistic threats are weaker yet structured. Label poisoning is a representative case in which workers flip a fraction of local labels but otherwise follow the protocol, producing poisoned messages through standard optimisation rather than arbitrary fabrication. Attacks of this type have been studied as data poisoning against FL systems \cite{10.1007/978-3-030-58951-6_24}, and dedicated defenses have been proposed that leverage properties of label flipping \cite{9897807}. Recent gradient-space defenses such as LFighter \cite{JEBREEL2024111} cluster worker gradients and discard small, dense clusters as potentially poisoned, but such clustering-based approaches tend to work best when worker data distributions are similar and can degrade under strong heterogeneity. A central recent finding is that robustness is regime-dependent: under label poisoning and sufficiently heterogeneous data, mean aggregation can be theoretically and empirically more robust than several robust aggregators \cite{ijcai2024p530}.

\section{Privacy and Deployment}
FL is frequently paired with privacy-enhancing primitives, yet model updates can still leak information; secure aggregation addresses this by allowing the server to recover only the sum of client updates without learning individual updates \cite{BonawitzIvanovKreuter:17}. This creates a deployability constraint for robustness: many robust aggregators and anomaly detectors rely on per-client visibility, which secure aggregation explicitly hides. One pragmatic direction is trust bootstrapping: FLTrust \cite{DBLP:journals/corr/abs-2012-13995} uses a server-side reference update direction to reweight client updates by alignment, improving robustness without requiring full pairwise comparisons across clients.

\section{Temporal Dynamics}
Temporal mechanisms such as server-side momentum and adaptive optimisation are widely used to stabilise training under noisy updates and partial participation \cite{reddi2021adaptive}, but temporal state also opens new attack and defense surfaces. An attacker may apply small, consistent perturbations across rounds, gradually biasing the server trajectory while remaining inconspicuous in any single round. History-aware robust optimisation leverages past information to improve resilience; learning-from-history frameworks \cite{pmlr-v139-karimireddy21a} construct robust directions using temporal structure. Attackers can also exploit history: HisMSA \cite{ZhangHisMSA:2025} models historical knowledge across rounds to craft stealthy poisoning updates that bypass multiple countermeasures, motivating defenses that track cross-round consistency rather than relying solely on round-wise outlier filtering.

\section{Lightweight History Statistics}
Scalable defenses must also be computationally efficient: high-dimensional updates make pairwise comparison, clustering, and repeated processing across rounds expensive, especially with large client populations and partial participation. A practical alternative is to track lightweight per-client summaries, such as update norms, residual scales, and residual-change statistics, rather than repeatedly computing full pairwise geometry. This type of state is cheap to maintain with exponential moving averages and is well matched to history-aware aggregation: the server can detect abrupt or persistent behavioural changes while avoiding a hard population-clustering assumption.

\section{Project Positioning}
Taken together, prior work suggests that robust FL is shaped by interacting factors: threat model strength and structure; aggregation design; hidden vulnerabilities under adaptive coordination; heterogeneity-driven robustness--utility trade-offs; privacy constraints; and temporal dynamics where both defenses and attacks exploit history. The recurring gaps are therefore (i) the variance shelter phenomenon enabling stealthy attacks under high benign variability \cite{BarBarGol:19}, (ii) misclassification of benign non-IID diversity as outliers by hard-filtering defenses \cite{ijcai2024p530,JEBREEL2024111}, (iii) limited observability under privacy mechanisms \cite{BonawitzIvanovKreuter:17}, and (iv) cross-round exploitation that defeats purely round-wise detectors \cite{ZhangHisMSA:2025}. This project is positioned to address these gaps by developing HRAC, a history-residual adaptive clipping strategy that combines robust global norm capping, per-client residual histories, and soft residual-change weighting.

\section{Conclusion}
In conclusion, this work investigates history-aware Byzantine robustness in federated learning through a soft-weighted aggregation design. The results highlight that widely used robust aggregators do not necessarily dominate simple averaging under all threat models and data heterogeneity regimes. This motivates a more scenario-driven view of robustness, where the choice of aggregation rule should be guided by the attack surface, the degree of non-IID drift, and the stability signals available to the server, rather than by worst-case robustness claims alone.

\end{bodytwocol}
